% Tabla estadistica CPU vs GPU para el proyecto GPU MD5 Cracker
% Fecha: 2026-05-25
% n = 5 repeticiones por escenario
\begin{table}[h]
\centering
\small
\begin{tabular}{l l c c c c c}
\hline
ID & Escenario & L & CPU (MH/s) & GPU (MH/s) & Speedup (MH/s) & Tiempo CPU/GPU \\
\hline
T1 & num encontrado & 6 & 6.2767 $\pm$ 0.4875 & 1307.38 $\pm$ 13.10 & 208.3x & 0.0200 / 0.0000 \\
T2 & num no encontrado (hash real) & 6 & 7.2685 $\pm$ 0.1439 & 1299.17 $\pm$ 9.86 & 178.7x & 0.1380 / 0.0000 \\
T3 & lower encontrado & 5 & 7.6567 $\pm$ 0.0255 & 1737.02 $\pm$ 9.16 & 226.9x & 0.8500 / 0.0100 \\
T4 & lower no encontrado (hash real) & 5 & 7.6756 $\pm$ 0.0353 & 1735.60 $\pm$ 10.36 & 226.1x & 1.5500 / 0.0100 \\
T5 & alnum encontrado & 6 & 7.3714 $\pm$ 0.0807 & 1835.01 $\pm$ 4.22 & 248.9x & 1.1500 / 30.9540 \\
T6 & alnum no encontrado (hash real) & 4 & 8.1346 $\pm$ 0.0462 & 1828.01 $\pm$ 2.93 & 224.7x & 1.8180 / 0.0100 \\
\hline
\end{tabular}
\caption{Comparacion estadistica de rendimiento entre implementaciones secuencial (CPU) y CUDA (GPU), con 5 repeticiones por escenario.}
\label{tab:comparacion_cpu_gpu_md5_stats}
\end{table}

% Notas tecnicas:
% 1) En T5, el tiempo total de GPU fue mayor pese al mayor throughput,
%    porque la version actual no detiene hilos dentro del batch al detectar coincidencia temprana.
% 2) En casos muy rapidos, el tiempo de GPU aparece como 0.0000 s por redondeo de impresion.
